\chapter{AGI capability hypotheses and limits}

Evolutor is an architecture and research program. It is not evidence of artificial general intelligence.

This chapter therefore replaces rhetorical AGI claims with a capability ledger:
\[
\boxed{\text{claim}=(capability,\ operational\ test,\ evidence,\ budget,\ limits)}.
\]

\section{Evidence ladder}
\Plate{01-evidence}{Every capability statement receives an evidence level.}{evo51:evidence}
Use six labels:
\begin{enumerate}
\item \textbf{implemented}: mechanism exists;
\item \textbf{tested}: local invariants/unit experiments pass;
\item \textbf{benchmark-demonstrated}: preregistered or held-out benchmark evidence exists;
\item \textbf{replicated}: independent or cross-lab reproduction exists;
\item \textbf{hypothesized}: architecture plausibly supports the capability but evidence is insufficient;
\item \textbf{speculative}: research direction without a validated mechanism.
\end{enumerate}
Implemented memory is not demonstrated continual learning; a planner is not demonstrated general reasoning.

\section{Breadth, depth and autonomy}
The Levels of AGI framework separates performance/depth from generality/breadth and treats autonomy as a deployment dimension. This is useful because a highly autonomous narrow system and a broad low-autonomy assistant pose different questions.

Evolutor records breadth, performance and autonomy separately rather than compressing them into one ``AGI score.''

\section{Capability taxonomy}
The research ledger tracks at least: generalization, transfer, reasoning, planning, continual learning, memory, adaptation, tool use, world modeling, self-improvement, metacognition, open-ended learning, autonomy and general intelligence.

\section{Generalization}
Operational question: does the system succeed on held-out distributions, tasks or environments not encoded during development? Report contamination controls, priors and adaptation budget.

In-distribution accuracy is useful skill evidence but insufficient evidence of general intelligence.

\section{Transfer}
Transfer asks whether learning on source experience improves acquisition or performance on a distinct target. Compare against a target-only baseline under matched budget.

\section{Reasoning}
Reasoning is not one scalar. Evaluate compositional, causal, spatial, symbolic, quantitative and interactive reasoning separately. Chain-of-thought-like text is not proof that the latent mechanism performed valid reasoning.

\section{Planning}
A planner must achieve goals over horizons under action costs, uncertainty and replanning. Measure success, actions, compute, violations and recovery. A static plan-shaped response is not enough.

\section{Continual learning}
Chapter 50 provides the protocol: acquisition plus retention/transfer under a stream. Updating durable memory alone is not parameter learning.

\section{Memory}
Memory has multiple forms: native engine state, active context, durable addressable records and learned parameters. Evaluate retention horizon, retrieval precision/recall, interference, provenance, deletion and security.

\section{Adaptation}
Adaptation includes test-time state change, retrieval, parameter/adaptor update, structural change and population search. Report which layer changed and its data/compute budget.

\section{Tool use}
Tool use is demonstrated when the system selects valid tools, forms correct typed calls, handles errors and completes tasks under deterministic authorization. Having an API connector is only implementation evidence.

\section{World modeling}
A world model should predict consequences or latent dynamics sufficiently to improve decisions in held-out environments. Narrative descriptions of the world are not necessarily predictive models.

\section{Metacognition}
Operationalize uncertainty calibration, error detection, strategy switching and knowing when to defer. Self-reported confidence without calibration is weak evidence.

\section{Self-improvement}
\Plate{02-self}{Improvement proposals remain subject to external evaluation and deployment authority.}{evo51:self}
A system that edits its code is not necessarily improving. Self-improvement requires repeated, causally attributed gains on held-out evaluations under fixed or declared budgets, without unacceptable regression.

\section{Open-ended learning}
Open-endedness requires sustained production/acquisition of new useful skills or niches rather than convergence on a fixed task set. It remains a research hypothesis for Evolutor.

\section{Autonomy}
Autonomy measures how much human intervention is required across task selection, planning, execution, recovery and continuation. Higher autonomy is not synonymous with higher intelligence and increases governance requirements.

\section{Skill-acquisition efficiency}
ARC-AGI emphasizes generalization and the efficiency with which new skills are acquired, not merely task-specific skill. ARC-AGI-2 adds explicit efficiency/cost reporting, and ARC-AGI-3 extends evaluation into interactive environments with exploration, goal acquisition, memory and action efficiency.

This motivates:
\[
\eta=\frac{\Delta capability}{data+compute+actions+risk},
\]
as a conceptual efficiency ratio. Its denominator is multi-dimensional; a single scalar requires a declared policy.

\section{Unknown-task evaluation}
\Plate{03-novel}{Unknown task/environment → limited experience → acquired skill.}{evo51:novel}
Holdouts must remain genuinely held out. Repeated manual tuning against a public evaluation set converts it into development data.

\section{Interactive evaluation}
Static benchmarks miss exploration, recovery and long-horizon interaction. Interactive environments can measure world-model acquisition, action efficiency, memory and planning, but they still cover only a bounded domain.

\section{DOGMA hypotheses}
DOGMA's compressed/native state motivates hypotheses about efficient persistent state and non-Transformer sequence processing. Those hypotheses require matched comparisons against GRU/SSM/Transformer and capacity/compute controls. Non-Transformer architecture is not itself evidence of greater generality.

\section{Hermon DNA hypotheses}
Hermon DNA preserves Transformer semantics while exploring DNA/genomic inductive biases and serving/runtime innovations. Better genomic benchmarks would establish domain skill, not AGI.

\section{Evolutor hypotheses}
Evolutor's compiler/runtime/memory/governance layers may support persistent, tool-using, adaptive systems. Integration breadth is an engineering property until held-out capability evaluations demonstrate generalization.

\section{Capability matrix}
\Plate{04-matrix}{Capability × benchmark × evidence × budget × failure matrix.}{evo51:matrix}
For every claimed capability record: definition, test suite, held-out policy, baseline, budget, metric, confidence interval, failure cases, evidence level and reproducibility status.

\section{No single benchmark}
ARC-like tests, software engineering, science, embodied/interactive tasks, continual learning and tool use probe different dimensions. Passing one does not entail the others.

\section{Safety is not intelligence}
\Plate{05-safety}{Capability evaluation and safety/governance evaluation are orthogonal axes.}{evo51:safety}
A capable system can be unsafe; a safe system can be limited. Evaluate authorization, isolation, monitoring, rollback and human control independently of capability scores.

\section{Stopping rules}
Research claims should have falsifiers. Example: if DOGMA does not outperform capacity/compute-matched recurrent/SSM baselines across preregistered long-context tasks and seeds, retract the claim that its state mechanism yields a general efficiency advantage.

\section{Executable claim ledger}
\Listing{reference}
The reference implements evidence levels, capability records and promotion rules. It deliberately prevents ``implemented'' from being promoted directly to ``demonstrated'' without benchmark evidence.

\section{Current conclusion}
The defensible statement is: Evolutor is a systems and research architecture for testing hypotheses about persistent, adaptive and governed AI. DOGMA and Hermon DNA are model/runtime research programs. None of the architecture, chapter or teaching-oracle results establish AGI.

\section{Twelve worked exercises}
\begin{enumerate}
\item Planner implemented means planning demonstrated? \textbf{No.}
\item Memory means continual learning? \textbf{No.}
\item Domain benchmark success means AGI? \textbf{No.}
\item Why held-out tasks? \textbf{To measure generalization rather than development fit.}
\item Autonomy equals intelligence? \textbf{No.}
\item Tool connector equals tool-use competence? \textbf{No.}
\item Why efficiency? \textbf{Brute-force skill can hide poor acquisition/generalization.}
\item Public eval repeatedly tuned against remains held out? \textbf{No.}
\item Self-editing equals self-improvement? \textbf{No.}
\item One benchmark sufficient? \textbf{No.}
\item Safety score part of intelligence score? \textbf{Track separately.}
\item Strongest current Evolutor claim? \textbf{Research architecture with testable hypotheses, not demonstrated AGI.}
\end{enumerate}

\section{Bridge}
Chapter 52 closes the book by consolidating failures, unresolved hypotheses, reproducibility artifacts and a falsifiable research agenda. The final synthesis should make it easier to disprove weak Evolutor claims, not harder.
